\chapter{高等技術}

本章探討更進階的 C 語言特性。 從命令列參數到函式指標，這些技術讓程式與操作系統和外部世界更好地互動。

\section{命令列參數解析}

程式可接收命令列參數。 \texttt{argc} 是參數個數，\texttt{argv} 是參數陣列。

\begin{code}{c}{command line args}
#include <stdio.h>

int main(int argc, char *argv[]) {
    printf("Program name: %s\n", argv[0]);
    printf("Number of arguments: %d\n", argc - 1);

    for (int i = 1; i < argc; i++) {
        printf("Argument %d: %s\n", i, argv[i]);
    }
    return 0;
}
\end{code}

執行：

\begin{code}{bash}{running with args}
$ ./program arg1 arg2 arg3
Program name: ./program
Number of arguments: 3
Argument 1: arg1
Argument 2: arg2
Argument 3: arg3
\end{code}

\keypoint{\texttt{argv[0]} 是程式名，\texttt{argv[1]} 到 \texttt{argv[argc-1]} 是引數。}

\section{命令列應用實例}

\subsection*{讀取並處理檔案}

\begin{code}{c}{file processing}
#include <stdio.h>
#include <stdlib.h>

int main(int argc, char *argv[]) {
    if (argc < 2) {
        fprintf(stderr, "Usage: %s <filename>\n", argv[0]);
        return 1;
    }

    FILE *fp = fopen(argv[1], "r");
    if (!fp) {
        perror("Cannot open file");
        return 1;
    }

    char line[256];
    while (fgets(line, sizeof(line), fp)) {
        printf("%s", line);
    }

    fclose(fp);
    return 0;
}
\end{code}

\subsection*{計算數字和}

\begin{example}{c}{(sum-args.c) Sum command line numbers}
#include <stdio.h>
#include <stdlib.h>

int main(int argc, char *argv[]) {
    if (argc < 2) {
        fprintf(stderr, "Usage: %s <num1> <num2> ...\n", argv[0]);
        return 1;
    }

    int sum = 0;
    for (int i = 1; i < argc; i++) {
        int num = atoi(argv[i]);
        sum += num;
    }

    printf("Sum: %d\n", sum);
    return 0;
}
\end{example}

\keypoint{命令列參數讓程式更靈活。 檔案名、設置都可透過參數傳遞。}

\section{全域變數與 extern}

全域變數跨檔案可見。 \texttt{extern} 宣告引用其他檔案的全域變數。

\subsection*{跨檔案全域變數}

檔案 1：\texttt{globals.c}

\begin{code}{c}{global variable definition}
int global_count = 0;

void increment() {
    global_count++;
}
\end{code}

檔案 2：\texttt{main.c}

\begin{code}{c}{extern variable}
#include <stdio.h>

extern int global_count;
void increment();

int main() {
    printf("Count: %d\n", global_count);
    increment();
    printf("Count: %d\n", global_count);
    return 0;
}
\end{code}

\subsection*{static 全域變數：檔案作用域}

\begin{code}{c}{static global}
// In globals.c
static int counter = 0;

void increment() {
    counter++;
}

int get_counter() {
    return counter;
}

// main.c cannot access counter directly
\end{code}

\keypoint{\texttt{extern} 引用外部全域變數。 \texttt{static} 限制作用域到單檔案。}

\section{函式指標}

指標可指向函式。 透過函式指標呼叫函式，實現回調和多態。

\begin{code}{c}{function pointers}
#include <stdio.h>

int add(int a, int b) { return a + b; }
int mul(int a, int b) { return a * b; }

int main() {
    int (*func_ptr)(int, int);

    func_ptr = add;
    printf("add: %d\n", func_ptr(3, 4));

    func_ptr = mul;
    printf("mul: %d\n", func_ptr(3, 4));

    return 0;
}
\end{code}

\subsection*{回調函式模式}

\begin{code}{c}{callback pattern}
#include <stdio.h>
#include <string.h>

typedef int (*CompareFunc)(const void*, const void*);

int compare_int(const void *a, const void *b) {
    return *(int*)a - *(int*)b;
}

void sort_with_callback(int *arr, int n,
                       CompareFunc compare) {
    // Simple bubble sort using callback
    for (int i = 0; i < n - 1; i++) {
        for (int j = 0; j < n - i - 1; j++) {
            if (compare(&arr[j], &arr[j+1]) > 0) {
                int temp = arr[j];
                arr[j] = arr[j+1];
                arr[j+1] = temp;
            }
        }
    }
}

int main() {
    int arr[] = {5, 2, 8, 1};
    sort_with_callback(arr, 4, compare_int);
    return 0;
}
\end{code}

\keypoint{函式指標實現回調。 \texttt{qsort} 和 \texttt{bsearch} 內部使用回調。}

\section{變長陣列 (VLA)}

C99 支援執行時決定的陣列大小。 但要小心棧空間限制。

\begin{code}{c}{variable length arrays}
#include <stdio.h>

int main() {
    int n;
    printf("Array size: ");
    scanf("%d", &n);

    int arr[n];  // Size determined at runtime

    for (int i = 0; i < n; i++) {
        arr[i] = i * i;
    }

    for (int i = 0; i < n; i++) {
        printf("%d ", arr[i]);
    }
    return 0;
}
\end{code}

\subsection*{VLA vs malloc}

VLA 快但有棧大小限制。 大陣列用 \texttt{malloc}。

\begin{code}{c}{VLA considerations}
// Small arrays: use VLA
int small[100];

// Large arrays: use malloc
int n = 1000000;
int *large = (int*)malloc(n * sizeof(int));
if (large) {
    // use large
    free(large);
}
\end{code}

\keypoint{VLA 方便但棧有限。 大於 MB 的陣列應用 \texttt{malloc}。}

\section{綜合應用}

多檔案程式使用上述技術實現模組化設計。

\subsection*{專案結構}

\begin{code}{text}{project layout}
project/
  Makefile
  math.h
  math.c
  utils.h
  utils.c
  main.c
\end{code}

\subsection*{header 檔案設計}

\texttt{math.h}：

\begin{code}{c}{header file}
#ifndef MATH_H
#define MATH_H

int add(int a, int b);
int sub(int a, int b);
int mul(int a, int b);

#endif
\end{code}

\subsection*{實作與連結}

\begin{code}{c}{implementation}
// math.c
#include "math.h"

int add(int a, int b) { return a + b; }
int sub(int a, int b) { return a - b; }
int mul(int a, int b) { return a * b; }

// main.c
#include <stdio.h>
#include "math.h"

int main() {
    printf("%d\n", add(5, 3));
    return 0;
}
\end{code}

編譯：

\begin{code}{bash}{makefile}
gcc -c math.c -o math.o
gcc -c main.c -o main.o
gcc math.o main.o -o program
\end{code}

\keypoint{多檔案程式靠 header 檔宣告、實作檔定義、Makefile 連結。}
